\section{Conclusion}
\label{sec:Conclusion}

In summary, it is far faster and easier to produce technical documents in \LaTeX\ once familiarity has been attained, and the results are of higher quality and appear more professional. Table~\ref{tab:Comparison} below gives a comparison of the key differences between Word and \LaTeX. The source code for this essay is available on my \href{https://github.com/HenryGinn/Essays/tree/main/TheCaseForLaTeX}{\underline{GitHub}}.

\begin{table}[H]
	\centering
	\caption{Comparing Word and \LaTeX.}
	\label{tab:Comparison}
	\begin{tabular}{ccc}
		\myhline
		Property & Word & \LaTeX  \\
		\myhline
		Editor format & What You See Is What You Get & Needs to compile  \\
		Bugfixing & Poor support and confusing behaviour & Good support with searchable errors  \\
		Software & Proprietary & Free and Open Source  \\
		Formatting control & Poor & High  \\
		Technical features & Lacking & Rich  \\
		Microsoft compatibility & High & Medium  \\
		Simplicity & Extremely simple &  $\sim$~1 hour of training  \\
		Version control & Medium & Git-compatible  \\
		\myhline
	\end{tabular}
\end{table}

\LaTeX\ is standard software in academia. Git is standard software in software engineering. They are both extremely mature and have been successfully implemented in these sectors so I see no reason why this could not also be the case in engineering and consultancy. Collaboration, version control, and producing technical documents is a hard problem, but it has been solved by \LaTeX\ and Git. Therefore I hope that WSP agrees that their use is inline with their digitial approach and considers implementing this technology.